\begingroup\small\setlength{\tabcolsep}{4pt}\renewcommand{\arraystretch}{1.16}
\begin{longtable}{lrp{.77\linewidth}}
\caption{Ví dụ lỗi từ AI consensus; chưa xác minh chuyên gia.}\label{tab:errors}\\
\toprule Query & Grade & Diễn giải AI \\
\midrule\endfirsthead
\multicolumn{3}{l}{\textit{Bảng tiếp theo}}\\
\toprule Query & Grade & Diễn giải AI \\
\midrule\endhead
\bottomrule\endfoot
Q11 & 0 & Tài liệu nghiên cứu về vật liệu cathode Na3V2(PO4)(2)O2F phủ rGO cho pin natri-ion, hoàn toàn không đề cập đến Ti3C2Tx MXene hay kỹ thuật chêm tầng (interlayer engineering) bằng MXene theo yêu cầu của truy vấn. \\
Q18 & 1 & Tài liệu đề cập đến phương pháp phục hồi (restoring) pin ion natri đã qua sử dụng thay vì thu hồi kim loại catốt (recovery of cathode metals) theo ý định nghiên cứu, do đó chỉ đạt mức độ nền tảng cơ bản hoặc một phần. \\
Q19 & 1 & Tài liệu đề cập đến pin ion nước (aqueous ion battery) nói chung, trong đó có nhắc đến pin natri-ion dạng nước (aqueous sodium-ion batteries). Tuy nhiên, dung dịch điện giải ở đây sử dụng dung môi hỗn hợp gồm nước và formamide (organic solvent is composed of organic solvent and water and the organic solvent is strongly polar formamide), chứ không phải là hệ điện giải muối đậm đặc 'water-in-salt electrolytes' như yêu cầu của truy vấn. Do đó, tài liệu chỉ đạt điểm 1 vì mang tính chất nền tảng rộng và thiếu điều kiện kỹ thuật cốt lõi. \\
\end{longtable}\endgroup
\source{outputs/retrieval_ai_evaluation_v1/ai_consensus_qrels.csv}
